\section{Konklusjon}
\label{sec:conclusion-nn}

Denne artikkelen spurde: kvifor tek eit designobjekt éi form framfor ei anna? Svaret frå \formlaere{} er at form er posisjonen i morforommet der formgrammatikken sin respons på tilpassingslandskapet sin gradient; slik agentane sine kognitive lyskjegler oppfattar han; produserte eit lokalt stabilt poly-agentisk kompromiss.

Teorien sameinar tre formelle tradisjonar: Stiny sin formalgebra gjev syntaksen til forma, Hatchuel og Weil sin C-K-teori gjev den epistemiske logikken, og Levin sitt multiskala-kompetanserammeverk gjev dynamikken. Syntesen produserer éi generativ likning: $\mathrm{Form} = SG(\nabla_{C(A)}\, L(c,t))$.

Teorien introduserer omsorg; dimensjonaliteten til agenten si kognitive lyskjegle; som mål på designintelligens. Ein agent som representerer fleire aksar i morforommet navigerer eit rikare landskap og produserer former som toler fleire trykk samstundes. Ein agent som snevrar lyskjegla si produserer former som sviktar langs dei aksane ho ikkje såg. Svikten er ikkje tilfeldig; ho er systematisk, og forma ber signaturen av blindskapen som produserte ho.

Me testa teorien mot 2\,066 europeiske stolar over 744 år. Tjue statistiske analysar over 14 kryssvalideringsdelmengder gjev resultat konsistente med kjernepredikasjonane i proposisjon 1--4; ingen postulat er falsifisert. Proposisjon 5--7 (agens, omsorg og den diagnostiske konklusjonen) er formaliserte og illustrerte, men krev direkte empirisk testing gjennom forskingsprogramma føreslått i seksjon~\ref{sec:agenda-nn}.

Me demonstrerte teoriens rekkevidde utover historisk møbeldesign gjennom fem utarbeidde døme: utkragingsstolkappløpet 1925--1928, Neri Oxman sitt Silk Pavilion, Figma Weave, samtaleagentdesign, og The Great Green Wall.

Teorien er falsifiserbar på kvart nivå. Kvar proposisjon genererer testbare predikasjonar. Sju umiddelbare forskingsprogram er føreslåtte.

Forma viser kor langt omsorga rakk. Teorien viser korleis me kan måle begge delar.
